% Options for packages loaded elsewhere
\PassOptionsToPackage{unicode}{hyperref}
\PassOptionsToPackage{hyphens}{url}
\PassOptionsToPackage{dvipsnames,svgnames,x11names}{xcolor}
%
\documentclass[
  11pt,
]{article}
\usepackage{amsmath,amssymb}
\usepackage{iftex}
\ifPDFTeX
  \usepackage[T1]{fontenc}
  \usepackage[utf8]{inputenc}
  \usepackage{textcomp} % provide euro and other symbols
\else % if luatex or xetex
  \usepackage{unicode-math} % this also loads fontspec
  \defaultfontfeatures{Scale=MatchLowercase}
  \defaultfontfeatures[\rmfamily]{Ligatures=TeX,Scale=1}
\fi
\usepackage{lmodern}
\ifPDFTeX\else
  % xetex/luatex font selection
\fi
% Use upquote if available, for straight quotes in verbatim environments
\IfFileExists{upquote.sty}{\usepackage{upquote}}{}
\IfFileExists{microtype.sty}{% use microtype if available
  \usepackage[]{microtype}
  \UseMicrotypeSet[protrusion]{basicmath} % disable protrusion for tt fonts
}{}
\makeatletter
\@ifundefined{KOMAClassName}{% if non-KOMA class
  \IfFileExists{parskip.sty}{%
    \usepackage{parskip}
  }{% else
    \setlength{\parindent}{0pt}
    \setlength{\parskip}{6pt plus 2pt minus 1pt}}
}{% if KOMA class
  \KOMAoptions{parskip=half}}
\makeatother
\usepackage{xcolor}
\usepackage[margin=1in]{geometry}
\usepackage{color}
\usepackage{fancyvrb}
\newcommand{\VerbBar}{|}
\newcommand{\VERB}{\Verb[commandchars=\\\{\}]}
\DefineVerbatimEnvironment{Highlighting}{Verbatim}{commandchars=\\\{\}}
% Add ',fontsize=\small' for more characters per line
\newenvironment{Shaded}{}{}
\newcommand{\AlertTok}[1]{\textcolor[rgb]{1.00,0.00,0.00}{\textbf{#1}}}
\newcommand{\AnnotationTok}[1]{\textcolor[rgb]{0.38,0.63,0.69}{\textbf{\textit{#1}}}}
\newcommand{\AttributeTok}[1]{\textcolor[rgb]{0.49,0.56,0.16}{#1}}
\newcommand{\BaseNTok}[1]{\textcolor[rgb]{0.25,0.63,0.44}{#1}}
\newcommand{\BuiltInTok}[1]{\textcolor[rgb]{0.00,0.50,0.00}{#1}}
\newcommand{\CharTok}[1]{\textcolor[rgb]{0.25,0.44,0.63}{#1}}
\newcommand{\CommentTok}[1]{\textcolor[rgb]{0.38,0.63,0.69}{\textit{#1}}}
\newcommand{\CommentVarTok}[1]{\textcolor[rgb]{0.38,0.63,0.69}{\textbf{\textit{#1}}}}
\newcommand{\ConstantTok}[1]{\textcolor[rgb]{0.53,0.00,0.00}{#1}}
\newcommand{\ControlFlowTok}[1]{\textcolor[rgb]{0.00,0.44,0.13}{\textbf{#1}}}
\newcommand{\DataTypeTok}[1]{\textcolor[rgb]{0.56,0.13,0.00}{#1}}
\newcommand{\DecValTok}[1]{\textcolor[rgb]{0.25,0.63,0.44}{#1}}
\newcommand{\DocumentationTok}[1]{\textcolor[rgb]{0.73,0.13,0.13}{\textit{#1}}}
\newcommand{\ErrorTok}[1]{\textcolor[rgb]{1.00,0.00,0.00}{\textbf{#1}}}
\newcommand{\ExtensionTok}[1]{#1}
\newcommand{\FloatTok}[1]{\textcolor[rgb]{0.25,0.63,0.44}{#1}}
\newcommand{\FunctionTok}[1]{\textcolor[rgb]{0.02,0.16,0.49}{#1}}
\newcommand{\ImportTok}[1]{\textcolor[rgb]{0.00,0.50,0.00}{\textbf{#1}}}
\newcommand{\InformationTok}[1]{\textcolor[rgb]{0.38,0.63,0.69}{\textbf{\textit{#1}}}}
\newcommand{\KeywordTok}[1]{\textcolor[rgb]{0.00,0.44,0.13}{\textbf{#1}}}
\newcommand{\NormalTok}[1]{#1}
\newcommand{\OperatorTok}[1]{\textcolor[rgb]{0.40,0.40,0.40}{#1}}
\newcommand{\OtherTok}[1]{\textcolor[rgb]{0.00,0.44,0.13}{#1}}
\newcommand{\PreprocessorTok}[1]{\textcolor[rgb]{0.74,0.48,0.00}{#1}}
\newcommand{\RegionMarkerTok}[1]{#1}
\newcommand{\SpecialCharTok}[1]{\textcolor[rgb]{0.25,0.44,0.63}{#1}}
\newcommand{\SpecialStringTok}[1]{\textcolor[rgb]{0.73,0.40,0.53}{#1}}
\newcommand{\StringTok}[1]{\textcolor[rgb]{0.25,0.44,0.63}{#1}}
\newcommand{\VariableTok}[1]{\textcolor[rgb]{0.10,0.09,0.49}{#1}}
\newcommand{\VerbatimStringTok}[1]{\textcolor[rgb]{0.25,0.44,0.63}{#1}}
\newcommand{\WarningTok}[1]{\textcolor[rgb]{0.38,0.63,0.69}{\textbf{\textit{#1}}}}
\usepackage{longtable,booktabs,array}
\usepackage{calc} % for calculating minipage widths
% Correct order of tables after \paragraph or \subparagraph
\usepackage{etoolbox}
\makeatletter
\patchcmd\longtable{\par}{\if@noskipsec\mbox{}\fi\par}{}{}
\makeatother
% Allow footnotes in longtable head/foot
\IfFileExists{footnotehyper.sty}{\usepackage{footnotehyper}}{\usepackage{footnote}}
\makesavenoteenv{longtable}
\setlength{\emergencystretch}{3em} % prevent overfull lines
\providecommand{\tightlist}{%
  \setlength{\itemsep}{0pt}\setlength{\parskip}{0pt}}
\setcounter{secnumdepth}{5}
\ifLuaTeX
  \usepackage{selnolig}  % disable illegal ligatures
\fi
\usepackage{bookmark}
\IfFileExists{xurl.sty}{\usepackage{xurl}}{} % add URL line breaks if available
\urlstyle{same}
\hypersetup{
  pdftitle={Petualangan Ajisaka - Technical Documentation},
  pdfauthor={Azzar},
  colorlinks=true,
  linkcolor={Maroon},
  filecolor={Maroon},
  citecolor={Blue},
  urlcolor={Blue},
  pdfcreator={LaTeX via pandoc}}

\title{Petualangan Ajisaka - Technical Documentation}
\author{Azzar}
\date{August 2026}

\begin{document}
\maketitle

{
\hypersetup{linkcolor=}
\setcounter{tocdepth}{3}
\tableofcontents
}
\section{Petualangan Ajisaka --- Aplikasi Pembelajaran Menulis Aksara
Jawa}\label{petualangan-ajisaka-aplikasi-pembelajaran-menulis-aksara-jawa}

\textbf{Design \& Development Document}

\begin{longtable}[]{@{}
  >{\raggedright\arraybackslash}p{(\columnwidth - 2\tabcolsep) * \real{0.5000}}
  >{\raggedright\arraybackslash}p{(\columnwidth - 2\tabcolsep) * \real{0.5000}}@{}}
\toprule\noalign{}
\endhead
\bottomrule\noalign{}
\endlastfoot
\textbf{Version} & 0.1 (draft) \\
\textbf{Date} & 2026-08-12 \\
\textbf{Status} & Design document --- pre-build \\
\textbf{Delivery target} & Installable Progressive Web App (PWA) \\
\textbf{Source material} & \texttt{projek\ Isif.pdf} --- \emph{Naskah
Lengkap Aplikasi Petualangan Ajisaka} \\
\end{longtable}

\begin{center}\rule{0.5\linewidth}{0.5pt}\end{center}

\subsection{1. Overview}\label{overview}

\textbf{Petualangan Ajisaka} is a gamified learning application that
teaches children to read and write \textbf{Aksara Jawa} (the Javanese
script) through a story-driven adventure. The player follows the
legendary hero \textbf{Ajisaka} across three islands, unlocking relics
and defeating enemies by completing on-screen handwriting practice
challenges.

The app runs entirely inside the browser and is packaged as an
installable, \textbf{offline-capable PWA} so students on low-end tablets
and school Chromebooks can play without connectivity.

\subsubsection{1.1 What the product is}\label{what-the-product-is}

\begin{itemize}
\tightlist
\item
  A \emph{stroke-practice} game: the player literally writes Aksara Jawa
  with a finger or stylus on a canvas, guided by animated arrow guides
  (stroke order + direction).
\item
  A \emph{narrative adventure}: writing missions are framed as story
  events (unlock a sword seal, answer a village elder's test, battle a
  giant's envoys).
\item
  A \emph{progression system}: three levels (Pemula → Mahir → Master),
  rewards (Pedang Pusaka, Perisai Sakti), followers joining the party
  (Dora, the local villager), and a final exam + crowned-ending.
\end{itemize}

\begin{center}\rule{0.5\linewidth}{0.5pt}\end{center}

\subsection{2. Goals \& Non-Goals}\label{goals-non-goals}

\subsubsection{2.1 Goals}\label{goals}

\begin{longtable}[]{@{}
  >{\raggedright\arraybackslash}p{(\columnwidth - 2\tabcolsep) * \real{0.3333}}
  >{\raggedright\arraybackslash}p{(\columnwidth - 2\tabcolsep) * \real{0.6667}}@{}}
\toprule\noalign{}
\begin{minipage}[b]{\linewidth}\raggedright
\#
\end{minipage} & \begin{minipage}[b]{\linewidth}\raggedright
Goal
\end{minipage} \\
\midrule\noalign{}
\endhead
\bottomrule\noalign{}
\endlastfoot
G1 & Teach correct \textbf{stroke order and direction} for Aksara Jawa
basics (Nglegena), Sandangan, and Pasangan. \\
G2 & Keep learners motivated through narrative, rewards, and escalating
difficulty. \\
G3 & Work \textbf{offline} and be \textbf{installable} on Android
tablets, iPads, and desktop browsers (PWA). \\
G4 & Give immediate, friendly feedback on each stroke (truthy/dirty,
live). \\
G5 & Track progress locally per device (no account required for MVP). \\
G6 & Provide a complete, reusable dataset of Aksara Jawa strokes and
Unicode mappings for future features (transliteration, quizzes). \\
\end{longtable}

\subsubsection{2.2 Non-goals (MVP)}\label{non-goals-mvp}

\begin{itemize}
\tightlist
\item
  ✗ Server accounts, cloud sync, or leaderboards (future work).
\item
  ✗ Freehand recognition beyond the taught question set (no arbitrary
  text input).
\item
  ✗ Audio instructions / speech synthesis (future; on-screen text only).
\item
  ✗ Multi-language UI beyond Indonesian + minimal English (future).
\item
  ✗ Native app stores (PWA only for now; TWA/APK wrapper considered
  later).
\end{itemize}

\begin{center}\rule{0.5\linewidth}{0.5pt}\end{center}

\subsection{3. Target Users \& Personas}\label{target-users-personas}

\begin{longtable}[]{@{}
  >{\raggedright\arraybackslash}p{(\columnwidth - 4\tabcolsep) * \real{0.3333}}
  >{\raggedright\arraybackslash}p{(\columnwidth - 4\tabcolsep) * \real{0.3333}}
  >{\raggedright\arraybackslash}p{(\columnwidth - 4\tabcolsep) * \real{0.3333}}@{}}
\toprule\noalign{}
\begin{minipage}[b]{\linewidth}\raggedright
Persona
\end{minipage} & \begin{minipage}[b]{\linewidth}\raggedright
Profile
\end{minipage} & \begin{minipage}[b]{\linewidth}\raggedright
Needs
\end{minipage} \\
\midrule\noalign{}
\endhead
\bottomrule\noalign{}
\endlastfoot
\textbf{Adi (9 y.o.)} & Primary: elementary student learning Aksara Jawa
in class & Playful, guided practice, immediate feedback, feels like a
game \\
\textbf{Bu Sari (32 y.o.)} & Teacher & Assignable practice, clear level
map, zero-install on classroom tablets \\
\textbf{Pak Damar (40 y.o.)} & Parent helping at home & Simple install,
offline, no account, safe kid UI \\
\end{longtable}

\textbf{Design consequence:} big touch targets (≥ 48px), high-contrast
solid buttons, readable Indonesian copy, minimal typing, zero ads, zero
tracking.

\begin{center}\rule{0.5\linewidth}{0.5pt}\end{center}

\subsection{4. Narrative \& Content Spec (from source
PDF)}\label{narrative-content-spec-from-source-pdf}

\subsubsection{4.1 High-level arc}\label{high-level-arc}

\begin{verbatim}
Petualangan Ajisaka
  ├─ Prolog ─ Mengenal Asal-usul Aksara Jawa (edukasi)
  ├─ Level 1 ─ Pemula   · Pulau Sanjaya    → Pedang Pusaka    (+ Dora joins)
  ├─ Level 2 ─ Mahir    · Pulau Adi Jaya   → Perisai Sakti    (+ warga lokal joins)
  └─ Level 3 ─ Master   · Kerajaan Nusantara (2 fase) → Raksasa Hijau defeated → crowned King
\end{verbatim}

\subsubsection{4.2 Screen / menu breakdown}\label{screen-menu-breakdown}

\begin{longtable}[]{@{}
  >{\raggedright\arraybackslash}p{(\columnwidth - 6\tabcolsep) * \real{0.0789}}
  >{\raggedright\arraybackslash}p{(\columnwidth - 6\tabcolsep) * \real{0.2105}}
  >{\raggedright\arraybackslash}p{(\columnwidth - 6\tabcolsep) * \real{0.2368}}
  >{\raggedright\arraybackslash}p{(\columnwidth - 6\tabcolsep) * \real{0.4737}}@{}}
\toprule\noalign{}
\begin{minipage}[b]{\linewidth}\raggedright
\#
\end{minipage} & \begin{minipage}[b]{\linewidth}\raggedright
Screen
\end{minipage} & \begin{minipage}[b]{\linewidth}\raggedright
Content
\end{minipage} & \begin{minipage}[b]{\linewidth}\raggedright
Writing mechanic
\end{minipage} \\
\midrule\noalign{}
\endhead
\bottomrule\noalign{}
\endlastfoot
0 & \textbf{Halaman Awal (Splash/Home)} & Media title + big
\textbf{PLAY} button & --- \\
1 & \textbf{Dashboard (Main Menu)} & 5 navigation menus: Prolog, Level
1, Level 2, Level 3, \emph{Kembali ke Rumah} & --- \\
2 & \textbf{Prolog} & Educational page: history \& origin of Aksara Jawa
& --- \\
3 & \textbf{Level 1 --- Pemula} (Pulau Sanjaya) & Sinopsis: unlock seal
+ take the sword. \textbf{Narasi sinopsis tidak diekstrak utuh dari PDF}
→ open question §20 & Write \textbf{Aksara Dasar (Nglegena)} with
\textbf{arrow stroke-guides} \\
4 & \textbf{Level 1 End} & Reward: \textbf{Pedang Pusaka}; meeting
\textbf{Dora} joins the party & --- \\
5 & \textbf{Level 2 --- Mahir} (Pulau Adi Jaya) & Sinopsis: sail with
Dora, find \textbf{Perisai Sakti}, a local villager sets a test &
\textbf{11 soal} (per narrative) of Sandangan writing \\
6 & \textbf{Level 2 End} & Reward: \textbf{Perisai Sakti}; villager
joins party; \emph{Next Level} & --- \\
7 & \textbf{Level 3 --- Master, Fase 1} (Penghadangan Dua Utusan) & Two
Green Giant envoys intercept the ship at sea; defeat them by writing &
\textbf{20 soal Aksara Pasangan} (stacked consonants: main letter +
attached pasangan below/side), full stroke-guide \\
8 & \textbf{Level 3 --- Master, Fase 2} (Penyegelan Raksasa Hijau) &
Infiltrate Nusantara Kingdom; free it from the Green Giant & \textbf{3
latihan menus × 5 soal = 15 soal}: (1) kalimat dgn Aksara Dasar, (2)
kalimat dgn Sandangan, (3) kalimat dgn Pasangan \\
9 & \textbf{Akhir Cerita} & Happy ending: Green Giant sealed away;
player crowned \textbf{Raja Kerajaan Nusantara} & --- \\
\end{longtable}

\subsubsection{4.3 Numbers locked from the
narrative}\label{numbers-locked-from-the-narrative}

\begin{longtable}[]{@{}lll@{}}
\toprule\noalign{}
Item & Count & Source \\
\midrule\noalign{}
\endhead
\bottomrule\noalign{}
\endlastfoot
Level 2 writing questions & \textbf{11} & ``ke-11 soal'' \\
Level 3 Fase 1 questions & \textbf{20} (pasangan) & narrative \\
Level 3 Fase 2 questions & \textbf{15} (3 × 5) & narrative \\
Level 1 question count & \emph{unspecified} & \textbf{open question
§20} \\
\end{longtable}

\subsubsection{4.4 Content gaps to confirm (from PDF
anomalies)}\label{content-gaps-to-confirm-from-pdf-anomalies}

\begin{enumerate}
\def\labelenumi{\arabic{enumi}.}
\tightlist
\item
  Level 1 synopsis text for Pulau Sanjaya (the ``unlock the seal'' beat
  is present, but the full story paragraph was empty in extraction).
\item
  Whether Level 1 \emph{requires} Nglegena only, or includes a tutorial
  sub-mode.
\item
  Reward naming convention + visual style (game-asset source or in-app
  SVG?).
\end{enumerate}

\begin{center}\rule{0.5\linewidth}{0.5pt}\end{center}

\subsection{5. Functional Requirements}\label{functional-requirements}

\subsubsection{FR-1 · Home \& Navigation}\label{fr-1-home-navigation}

\begin{itemize}
\tightlist
\item
  \textbf{FR-1.1} Home screen shows the media title and a primary
  \textbf{PLAY} button.
\item
  \textbf{FR-1.2} PLAY opens the dashboard; system back or ``Kembali ke
  Rumah'' returns home.
\item
  \textbf{FR-1.3} Dashboard always lists \texttt{Prolog},
  \texttt{Level\ 1}, \texttt{Level\ 2}, \texttt{Level\ 3},
  \texttt{Kembali\ ke\ Rumah}.
\item
  \textbf{FR-1.4} Locked levels are visually locked and show a
  ``complete previous level'' hint (do not hard-block navigation; allow
  reading Prolog anytime).
\end{itemize}

\subsubsection{FR-2 · Story / Prolog}\label{fr-2-story-prolog}

\begin{itemize}
\tightlist
\item
  \textbf{FR-2.1} Prolog renders the origin story of Aksara Jawa as
  illustrated scenes.
\item
  \textbf{FR-2.2} Supports tap-to-advance (like a slide deck); skippable
  with confirmation.
\end{itemize}

\subsubsection{FR-3 · Writing Practice (core
gameplay)}\label{fr-3-writing-practice-core-gameplay}

\begin{itemize}
\tightlist
\item
  \textbf{FR-3.1} Each question shows a reference glyph (rendered with
  the Javanese font) + its romanization + Indonesian hint.
\item
  \textbf{FR-3.2} A canvas area captures freehand strokes (single
  touch/stylus; multi-touch disabled).
\item
  \textbf{FR-3.3} Animated \textbf{stroke guide}: numbered arrow heads
  showing direction and order, replayable on demand.
\item
  \textbf{FR-3.4} Live feedback: stroke match indicator (e.g., green =
  good, amber = acceptable, red = wrong direction/order).
\item
  \textbf{FR-3.5} Between strokes, guide for the next sub-stroke
  highlights; completed strokes stay visible.
\item
  \textbf{FR-3.6} \textbf{Undo / Clear} controls (big, kid-safe).
\item
  \textbf{FR-3.7} After a question completes: celebratory
  micro-feedback, then auto-advance or ``Lanjut''.
\item
  \textbf{FR-3.8} A wrong-but-attempted answer can be retried (no
  fail-state lock on MVP).
\end{itemize}

\subsubsection{FR-4 · Level Flow \&
Rewards}\label{fr-4-level-flow-rewards}

\begin{itemize}
\tightlist
\item
  \textbf{FR-4.1} Level intro screen = story beat (sinopsis) → practice
  session → outro beat.
\item
  \textbf{FR-4.2} On completion show reward item + story update (e.g.,
  ``Dora bergabung!'').
\item
  \textbf{FR-4.3} ``Next Level'' button unlocks the following level.
\item
  \textbf{FR-4.4} Level 3 has two phases with a phase-transition
  interstitial.
\end{itemize}

\subsubsection{FR-5 · Persistence}\label{fr-5-persistence}

\begin{itemize}
\tightlist
\item
  \textbf{FR-5.1} Persist: completed levels, per-question best score,
  collected rewards, settings.
\item
  \textbf{FR-5.2} Resumable: closing mid-level returns to dashboard with
  progress intact; continue begins from first unanswered question.
\item
  \textbf{FR-5.3} Storage: local-first (IndexedDB via
  \texttt{localStorage} facade for MVP state; canvas data ephemeral).
\end{itemize}

\subsubsection{FR-6 · PWA}\label{fr-6-pwa}

\begin{itemize}
\tightlist
\item
  \textbf{FR-6.1} Installable (valid manifest + icons + service worker).
\item
  \textbf{FR-6.2} Official \textbf{offline-first}: full app shell +
  asset (fonts, data, icons) precache; no network needed to play.
\item
  \textbf{FR-6.3} Updates: versioned SW with prompt-to-refresh when a
  new release is cached.
\end{itemize}

\begin{center}\rule{0.5\linewidth}{0.5pt}\end{center}

\subsection{6. Information Architecture \&
Navigation}\label{information-architecture-navigation}

\begin{verbatim}
Home (Halaman Awal)
  │  [PLAY]
  ▼
Dashboard ────────────────► (5 menu)
  ├── Prolog ─────────────► Story slides ──► back
  ├── Level 1 (Pemula) ───► [Intro sinopsis] ─► Practice L1 ─► [Reward: Pedang] ─► Next Level
  ├── Level 2 (Mahir) ────► [Intro sinopsis] ─► Practice L2 (11) ─► [Reward: Perisai] ─► Next Level
  ├── Level 3 (Master) ───► Fase 1: Practice Pasangan (20) ─► Fase 2: Lat.1 (5) Lat.2 (5) Lat.3 (5) ─► Ending
  └── Kembali ke Rumah ───► Home
\end{verbatim}

\subsubsection{6.1 Navigational rules}\label{navigational-rules}

\begin{itemize}
\tightlist
\item
  Global, persistent \textbf{back} affordance on every non-dashboard
  screen (system back + visible button).
\item
  Dashboard is the central hub; all level screens return to it.
\item
  \emph{Next Level} is the only forward gate; it appears only after a
  level completes.
\end{itemize}

\begin{center}\rule{0.5\linewidth}{0.5pt}\end{center}

\subsection{7. User Flows}\label{user-flows}

\subsubsection{7.1 First play through Level
1}\label{first-play-through-level-1}

\begin{Shaded}
\begin{Highlighting}[]
\NormalTok{flowchart TD}
\NormalTok{  A[Home] {-}{-}\textgreater{}|PLAY| B[Dashboard]}
\NormalTok{  B {-}{-}\textgreater{}|Level 1| C["Intro sinopsis: Pulau Sanjaya"]}
\NormalTok{  C {-}{-}\textgreater{} D["Writing practice: Nglegena set"]}
\NormalTok{  D {-}{-}\textgreater{}|question done| D}
\NormalTok{  D {-}{-}\textgreater{}|all done| E["Outro: Pedang Pusaka + Dora joins"]}
\NormalTok{  E {-}{-}\textgreater{}|Next Level| F[Dashboard {-} Level 2 unlocked]}
\end{Highlighting}
\end{Shaded}

\subsubsection{7.2 Single-question writing
loop}\label{single-question-writing-loop}

\begin{Shaded}
\begin{Highlighting}[]
\NormalTok{flowchart TD}
\NormalTok{  A[Show glyph + hint] {-}{-}\textgreater{} B[Replay stroke guide]}
\NormalTok{  B {-}{-}\textgreater{} C[Player traces on canvas]}
\NormalTok{  C {-}{-}\textgreater{} D\{Stroke match engine\}}
\NormalTok{  D {-}{-}\textgreater{}|pass| E["Stroke locked green {-} next guide"]}
\NormalTok{  D {-}{-}\textgreater{}|dirty {-} retry stroke| C}
\NormalTok{  E {-}{-}\textgreater{} F\{All sub{-}strokes done?\}}
\NormalTok{  F {-}{-}\textgreater{}|no| B}
\NormalTok{  F {-}{-}\textgreater{}|yes| G["Celebration {-} record score"]}
\NormalTok{  G {-}{-}\textgreater{} H\{More questions?\}}
\NormalTok{  H {-}{-}\textgreater{}|yes| A}
\NormalTok{  H {-}{-}\textgreater{}|no| I[Level outro + reward]}
\end{Highlighting}
\end{Shaded}

\begin{center}\rule{0.5\linewidth}{0.5pt}\end{center}

\subsection{8. Screen Inventory (route
map)}\label{screen-inventory-route-map}

\begin{longtable}[]{@{}
  >{\raggedright\arraybackslash}p{(\columnwidth - 4\tabcolsep) * \real{0.3333}}
  >{\raggedright\arraybackslash}p{(\columnwidth - 4\tabcolsep) * \real{0.3333}}
  >{\raggedright\arraybackslash}p{(\columnwidth - 4\tabcolsep) * \real{0.3333}}@{}}
\toprule\noalign{}
\begin{minipage}[b]{\linewidth}\raggedright
Route
\end{minipage} & \begin{minipage}[b]{\linewidth}\raggedright
Screen
\end{minipage} & \begin{minipage}[b]{\linewidth}\raggedright
Key UI
\end{minipage} \\
\midrule\noalign{}
\endhead
\bottomrule\noalign{}
\endlastfoot
\texttt{/} & Home & Title, PLAY \\
\texttt{/menu} & Dashboard & 5 level cards \\
\texttt{/prolog} & Prolog story & Slides, skip \\
\texttt{/level/1} & L1 flow & Intro / practice / outro \\
\texttt{/level/2} & L2 flow & Intro / practice(11) / outro \\
\texttt{/level/3} & L3 flow & Fase1 practice(20) → Fase2 3×latihan(15) →
ending \\
\texttt{/ending} & Happy ending & Crown ceremony \\
\end{longtable}

(Routes are client-side only; PWA serves the app shell for every path →
use hash router or SW navigation fallback to \texttt{/}.)

\begin{center}\rule{0.5\linewidth}{0.5pt}\end{center}

\subsection{9. Design System}\label{design-system}

\subsubsection{9.1 Design language}\label{design-language}

Three layers (per product guardrails):

\begin{enumerate}
\def\labelenumi{\arabic{enumi}.}
\tightlist
\item
  \textbf{Material You} --- interaction model: adaptive touch surfaces,
  tonal surfaces, dynamic elevation, large padded buttons, rounded
  corners.
\item
  \textbf{Minimalism} --- layout discipline: generous whitespace, clear
  hierarchy, one primary action per screen, reduced visual noise so a
  9-year-old isn't overwhelmed.
\item
  \textbf{Glassmorphism} --- only as an \emph{accent} on
  story/interstitial screens (frosted panels over illustrated
  backgrounds); not on gameplay surfaces (sketch canvases must stay
  glare-free for writing).
\end{enumerate}

Genre: \textbf{playful} (post-Linear soft school: consumer, casual,
children). Warm, pastel-dominant palette with deep ``keraton'' accents.
Illustrations are the brand: hand-drawn Ajisaka line-art silhouettes
(simple SVG), not photos.

\subsubsection{9.2 Palette (Keraton theme)}\label{palette-keraton-theme}

Derived from the playful \textbf{Carnival} token system, re-hued toward
Javanese heritage colors (keraton indigo/red-gold + batik neutrals),
pastel-first.

\begin{longtable}[]{@{}
  >{\raggedright\arraybackslash}p{(\columnwidth - 4\tabcolsep) * \real{0.3333}}
  >{\raggedright\arraybackslash}p{(\columnwidth - 4\tabcolsep) * \real{0.3333}}
  >{\raggedright\arraybackslash}p{(\columnwidth - 4\tabcolsep) * \real{0.3333}}@{}}
\toprule\noalign{}
\begin{minipage}[b]{\linewidth}\raggedright
Token
\end{minipage} & \begin{minipage}[b]{\linewidth}\raggedright
Value (OKLCH)
\end{minipage} & \begin{minipage}[b]{\linewidth}\raggedright
Use
\end{minipage} \\
\midrule\noalign{}
\endhead
\bottomrule\noalign{}
\endlastfoot
\texttt{-\/-color-paper} & \texttt{oklch(97\%\ 0.012\ 78)} & App
background (batik cream) \\
\texttt{-\/-color-paper-2} & \texttt{oklch(92\%\ 0.03\ 78)} & Card
surfaces \\
\texttt{-\/-color-paper-3} & \texttt{oklch(84\%\ 0.05\ 55)} &
Hover/interactive surface \\
\texttt{-\/-color-text} & \texttt{oklch(34\%\ 0.04\ 262)} & Primary text
(deep indigo-black) \\
\texttt{-\/-color-text-2} & \texttt{oklch(52\%\ 0.03\ 262)} & Muted
text \\
\texttt{-\/-color-accent} & \texttt{oklch(58\%\ 0.14\ 25)} & Primary
action (keraton red / terracotta) \\
\texttt{-\/-color-accent-2} & \texttt{oklch(70\%\ 0.13\ 80)} & Success /
gold (pedang \& perisai rewards) \\
\texttt{-\/-color-warn} & \texttt{oklch(66\%\ 0.14\ 60)} & Warning /
retry stroke \\
\texttt{-\/-color-error} & \texttt{oklch(52\%\ 0.16\ 27)} & Error /
wrong stroke \\
\texttt{-\/-color-border} & \texttt{oklch(86\%\ 0.025\ 78)} &
Dividers \\
\texttt{-\/-color-focus} & \texttt{oklch(60\%\ 0.16\ 262)} &
\texttt{:focus-visible} ring (indigo) \\
\texttt{-\/-color-glass} & \texttt{oklch(100\%\ 0\ 0\ /\ 0.55)} &
Frosted overlay fill \\
\end{longtable}

Dark mode: derived automatically by inverting paper band (keep for
classroom) via \texttt{@media\ (prefers-color-scheme:\ dark)} +
\texttt{{[}data-theme=dark{]}}.

\subsubsection{9.3 Typography}\label{typography}

\begin{longtable}[]{@{}
  >{\raggedright\arraybackslash}p{(\columnwidth - 4\tabcolsep) * \real{0.3333}}
  >{\raggedright\arraybackslash}p{(\columnwidth - 4\tabcolsep) * \real{0.3333}}
  >{\raggedright\arraybackslash}p{(\columnwidth - 4\tabcolsep) * \real{0.3333}}@{}}
\toprule\noalign{}
\begin{minipage}[b]{\linewidth}\raggedright
Role
\end{minipage} & \begin{minipage}[b]{\linewidth}\raggedright
Token
\end{minipage} & \begin{minipage}[b]{\linewidth}\raggedright
Stack
\end{minipage} \\
\midrule\noalign{}
\endhead
\bottomrule\noalign{}
\endlastfoot
Display / game titles & \texttt{-\/-font-display} &
\texttt{\textquotesingle{}Fredoka\ One\textquotesingle{},\ \textquotesingle{}Baloo\ 2\textquotesingle{},\ system-ui,\ sans-serif}
(rounded, kid-friendly) \\
Body / copy & \texttt{-\/-font-body} &
\texttt{\textquotesingle{}Nunito\textquotesingle{},\ system-ui,\ sans-serif} \\
Code / glyph data labels & \texttt{-\/-font-mono} &
\texttt{\textquotesingle{}JetBrains\ Mono\textquotesingle{},\ monospace} \\
\end{longtable}

\textbf{Scale} (from token system): display
\texttt{clamp(2.5rem,6vw,4.5rem)}, display-s
\texttt{clamp(2rem,4vw,3rem)}, 2xl \texttt{1.75rem}, xl
\texttt{1.25rem}, lg \texttt{1.125rem}, base \texttt{1rem}, sm
\texttt{0.875rem}, xs \texttt{0.75rem}.

\textbf{Rules:} headings always roman (no italic headings); italics only
inside body copy for emphasis. Indonesian copy throughout; keep
sentences short (≤ 12 words).

\textbf{Javanese glyph font:} self-hosted \textbf{Noto Sans Javanese}
(Unicode block U+A980--U+A9DF) for rendering reference glyphs. Fallback
font-stack on canvas:
\texttt{\textquotesingle{}Noto\ Sans\ Javanese\textquotesingle{},\ sans-serif}.
(Earliest glyph loads are precached by the SW.)

\subsubsection{9.4 Spacing, radius,
elevation}\label{spacing-radius-elevation}

\begin{itemize}
\tightlist
\item
  Spacing: 4px base scale (xs 4, sm 8, md 12, lg 16, xl 24, 2xl 32, 3xl
  48 \ldots).
\item
  Radius: sm 4 / md 8 / lg 16 --- reuse, never invent.
\item
  Elevation: Material tonal elevation --- cards = paper-2 on paper;
  modals use \texttt{box-shadow} with glass blur (story interstitials
  only).
\item
  Touch targets: \textbf{min 48×48px}, interactive gap ≥ 8px.
\end{itemize}

\subsubsection{9.5 Motion}\label{motion}

\begin{itemize}
\tightlist
\item
  Durations: fast 150ms (hover/micro), base 250ms, slow 400ms (screen
  transitions).
\item
  Easings: \texttt{cubic-bezier(0.16,1,0.3,1)} enter /
  \texttt{cubic-bezier(0.4,0,0.68,0.06)} exit.
\item
  All entrance/stagger/counter/loader animation through \textbf{anime.js
  v4} snippets, each guarded by \texttt{prefers-reduced-motion}.
\item
  Story slides: crossfade + slight scale; gameplay feedback: quick
  spring on stroke-lock.
\item
  \textbf{Game feel:} confetti-ish particle burst on question completion
  and at level rewards (Web Canvas, keep under 60 particles).
\end{itemize}

\subsubsection{9.6 Iconography}\label{iconography}

Font Awesome (free) per defaults; considered alternatives: Material
Symbols, Lucide. Use spellings/visuals a child recognizes (shield,
sword, home, arrow, question, star). All icon buttons carry
\texttt{aria-label}.

\subsubsection{9.7 Components (8-state
discipline)}\label{components-8-state-discipline}

\begin{longtable}[]{@{}
  >{\raggedright\arraybackslash}p{(\columnwidth - 2\tabcolsep) * \real{0.5000}}
  >{\raggedright\arraybackslash}p{(\columnwidth - 2\tabcolsep) * \real{0.5000}}@{}}
\toprule\noalign{}
\begin{minipage}[b]{\linewidth}\raggedright
Component
\end{minipage} & \begin{minipage}[b]{\linewidth}\raggedright
States to style
\end{minipage} \\
\midrule\noalign{}
\endhead
\bottomrule\noalign{}
\endlastfoot
\texttt{Button} (primary/ghost/danger) & default, hover, focus-visible,
active, disabled, loading, error, success \\
\texttt{LevelCard} & default, hover, focus, active, locked, completed,
current, disabled \\
\texttt{PracticeCanvas} & idle, guiding, drawing, stroke-locked,
feedback(ok/warn/error), cleared \\
\texttt{IconButton} (back/undo/clear/replay) & all 8 states \\
\texttt{Modal} (congrats, reset, update) & open/close anim + focus
trap \\
\end{longtable}

\begin{center}\rule{0.5\linewidth}{0.5pt}\end{center}

\subsection{10. Domain Model \& Data}\label{domain-model-data}

\subsubsection{\texorpdfstring{10.1 Aksara dataset
(\texttt{src/data/aksara.ts})}{10.1 Aksara dataset (src/data/aksara.ts)}}\label{aksara-dataset-srcdataaksara.ts}

\begin{Shaded}
\begin{Highlighting}[]
\KeywordTok{type}\NormalTok{ AksaraType }\OperatorTok{=} \StringTok{\textquotesingle{}nglegena\textquotesingle{}} \OperatorTok{|} \StringTok{\textquotesingle{}sandangan\textquotesingle{}} \OperatorTok{|} \StringTok{\textquotesingle{}pasangan\textquotesingle{}}\OperatorTok{;}

\KeywordTok{interface}\NormalTok{ AksaraGlyph \{}
\NormalTok{  id}\OperatorTok{:} \DataTypeTok{string}\OperatorTok{;}                    \CommentTok{// "ha", "na", "ca" ...}
\NormalTok{  type}\OperatorTok{:}\NormalTok{ AksaraType}\OperatorTok{;}
\NormalTok{  label}\OperatorTok{:} \DataTypeTok{string}\OperatorTok{;}                 \CommentTok{// romanization}
\NormalTok{  hintID}\OperatorTok{:} \DataTypeTok{string}\OperatorTok{;}                \CommentTok{// Indonesian hint translation key}
\NormalTok{  unicode}\OperatorTok{:} \DataTypeTok{string}\OperatorTok{;}               \CommentTok{// e.g. "\textbackslash{}uA98F" (ha)}
\NormalTok{  unicodePasangan}\OperatorTok{?:} \DataTypeTok{string}\OperatorTok{;}      \CommentTok{// pasangan glyph codepoint where applicable}
  \CommentTok{/** Normalized vector strokes, in units of the write{-}box (0..1). */}
\NormalTok{  strokes}\OperatorTok{:}\NormalTok{ Stroke[]}\OperatorTok{;}             \CommentTok{// ordered! = stroke{-}order/isomorphism}
\NormalTok{  sound}\OperatorTok{?:} \DataTypeTok{string}\OperatorTok{;}                \CommentTok{// phoneme for future TTS}
\NormalTok{\}}

\KeywordTok{interface}\NormalTok{ Stroke \{}
\NormalTok{  points}\OperatorTok{:} \BuiltInTok{Array}\OperatorTok{\textless{}}\NormalTok{\{ x}\OperatorTok{:} \DataTypeTok{number}\OperatorTok{;}\NormalTok{ y}\OperatorTok{:} \DataTypeTok{number}\NormalTok{ \}}\OperatorTok{\textgreater{};}  \CommentTok{// polyline, 0..1 box coords}
\NormalTok{  tolerance}\OperatorTok{:} \DataTypeTok{number}\OperatorTok{;}                         \CommentTok{// per{-}stroke looseness (default 0.12)}
\NormalTok{\}}
\end{Highlighting}
\end{Shaded}

\subsubsection{10.2 Level config}\label{level-config}

\begin{Shaded}
\begin{Highlighting}[]
\KeywordTok{interface}\NormalTok{ LevelConfig \{}
\NormalTok{  id}\OperatorTok{:} \DataTypeTok{number}\OperatorTok{;}
\NormalTok{  title}\OperatorTok{:} \DataTypeTok{string}\OperatorTok{;}          \CommentTok{// "Level 1 · Pemula · Pulau Sanjaya"}
\NormalTok{  questions}\OperatorTok{:}\NormalTok{ Question[]}\OperatorTok{;}
\NormalTok{  reward}\OperatorTok{:}\NormalTok{ \{ id}\OperatorTok{:} \DataTypeTok{string}\OperatorTok{;}\NormalTok{ name}\OperatorTok{:} \DataTypeTok{string}\NormalTok{ \}}\OperatorTok{;}   \CommentTok{// "Pedang Pusaka"}
\NormalTok{  story}\OperatorTok{:}\NormalTok{ \{ intro}\OperatorTok{:} \DataTypeTok{string}\OperatorTok{;}\NormalTok{ outro}\OperatorTok{:} \DataTypeTok{string}\NormalTok{ \}}\OperatorTok{;}
\NormalTok{\}}

\KeywordTok{interface}\NormalTok{ Question \{}
\NormalTok{  id}\OperatorTok{:} \DataTypeTok{string}\OperatorTok{;}
\NormalTok{  glyphId}\OperatorTok{:} \DataTypeTok{string}\OperatorTok{;}        \CommentTok{// → AksaraGlyph}
\NormalTok{  prompt}\OperatorTok{:} \DataTypeTok{string}\OperatorTok{;}         \CommentTok{// "Tulis aksara: HA (dasar)"}
\NormalTok{  sentence}\OperatorTok{?:} \DataTypeTok{string}\OperatorTok{;}      \CommentTok{// full Javanese sentence for Fase{-}2 questions}
\NormalTok{\}}
\end{Highlighting}
\end{Shaded}

\subsubsection{10.3 Seed content (proposed; volumes = open
questions)}\label{seed-content-proposed-volumes-open-questions}

\begin{longtable}[]{@{}
  >{\raggedright\arraybackslash}p{(\columnwidth - 4\tabcolsep) * \real{0.3333}}
  >{\raggedright\arraybackslash}p{(\columnwidth - 4\tabcolsep) * \real{0.3333}}
  >{\raggedright\arraybackslash}p{(\columnwidth - 4\tabcolsep) * \real{0.3333}}@{}}
\toprule\noalign{}
\begin{minipage}[b]{\linewidth}\raggedright
Set
\end{minipage} & \begin{minipage}[b]{\linewidth}\raggedright
Type
\end{minipage} & \begin{minipage}[b]{\linewidth}\raggedright
Count
\end{minipage} \\
\midrule\noalign{}
\endhead
\bottomrule\noalign{}
\endlastfoot
Nglegena (dasar) & 20 + optional variants & used in L1, L3-Fase2 menu
1 \\
Sandangan & 8 core (a→i,u,e,o, taling/tarung, pepet, cecak, etc.) & L2
(11 soal incl.~combos), L3-F2 menu 2 \\
Pasangan & pair of main + pasangan glyph & L3-F1 (20 soal), L3-F2 menu
3 \\
\end{longtable}

Final counts per level confirmed in §4.3; question contents (which
glyphs) are a \textbf{curriculum decision} --- propose defaults §20.

\begin{center}\rule{0.5\linewidth}{0.5pt}\end{center}

\subsection{11. Stroke Engine Design (core
challenge)}\label{stroke-engine-design-core-challenge}

\subsubsection{11.1 Capture}\label{capture}

\begin{itemize}
\tightlist
\item
  \texttt{\textless{}canvas\textgreater{}} full viewport of the
  write-box; \textbf{Pointer Events} (\texttt{pointerdown/move/up})
  unified for touch + stylus + mouse; \texttt{touch-action:\ none} to
  stop scrolling.
\item
  Points downsampled (≥ 10px spacing) and normalized into the 0..1
  write-box.
\item
  HiDPI: canvas sized to \texttt{devicePixelRatio}, CSS box fixed.
\end{itemize}

\subsubsection{11.2 Stroke-guide
rendering}\label{stroke-guide-rendering}

\begin{itemize}
\tightlist
\item
  Draw each reference \texttt{Stroke} polyline as faint target guide.
\item
  Animated progress along the path = the classic ``arrow head'' tracer
  (anime.js or requestAnimationFrame), replayable via \textbf{Replay}
  button.
\item
  Next expected sub-stroke highlighted; completed strokes rendered solid
  + check.
\end{itemize}

\subsubsection{11.3 Online matching (offline, deterministic,
fast)}\label{online-matching-offline-deterministic-fast}

Approach to keep 0 weight and run on-device: 1. \textbf{Normalize} the
input stroke (resample to N=32 points, scale to 0..1, translate to
center, keep proportional aspect). 2. \textbf{Direction signature:}
per-segment angle buckets (e.g., 8 bins). 3. \textbf{Dynamic Time
Warping (DTW)} distance to the reference normalized stroke (proxy: use
multi-step/rough DTW on resampled polylines). 4. \textbf{Order = truth:}
strokes must be completed in the reference order; a reversal triggers
the ``wrong direction'' feedback even if the shape matches. 5. Score →
\texttt{pass\ /\ warn\ /\ retry} via per-stroke tolerance.

Optional future upgrade: a tiny \textbf{ONNX/TFLite} model
(\texttt{12\ class\ output}, \textless{} 500 KB) trained on
augmentations of the 32 reference glyphs, run via
\texttt{@xenova/transformers} or plain TensorFlow.js --- swapped behind
the same \texttt{StrokeMatcher} interface. \textbf{Not in MVP.}

\subsubsection{11.4 Feedback mapping}\label{feedback-mapping}

\begin{longtable}[]{@{}lll@{}}
\toprule\noalign{}
Signal & Meaning & Color \\
\midrule\noalign{}
\endhead
\bottomrule\noalign{}
\endlastfoot
pass & stroke matched shape+order & green (\texttt{accent-2}) \\
warn & shape ok, slight drift & gold (\texttt{warn}) \\
retry & wrong direction / order & red (\texttt{error}) \\
\end{longtable}

\begin{center}\rule{0.5\linewidth}{0.5pt}\end{center}

\subsection{12. PWA \& Architecture}\label{pwa-architecture}

\subsubsection{12.1 Recommended stack}\label{recommended-stack}

\begin{longtable}[]{@{}
  >{\raggedright\arraybackslash}p{(\columnwidth - 2\tabcolsep) * \real{0.5000}}
  >{\raggedright\arraybackslash}p{(\columnwidth - 2\tabcolsep) * \real{0.5000}}@{}}
\toprule\noalign{}
\begin{minipage}[b]{\linewidth}\raggedright
Concern
\end{minipage} & \begin{minipage}[b]{\linewidth}\raggedright
Choice
\end{minipage} \\
\midrule\noalign{}
\endhead
\bottomrule\noalign{}
\endlastfoot
Language & \textbf{TypeScript (strict)} \\
Build & \textbf{Vite} \\
UI framework & \textbf{React 18} (or Svelte if a lighter runtime is
preferred --- see §20) \\
Styling & \textbf{Tailwind CSS v4} with CSS custom-property tokens
(OKLCH) \\
State & Zustand with \texttt{persist} middleware (localStorage) \\
Routing & \texttt{react-router} + \textbf{hash routing}
(offline-friendly, no server rewrites needed) \\
PWA & \texttt{vite-plugin-pwa} (Workbox): manifest + SW precache +
offline fallback \\
Canvas/game & Native Canvas 2D + anime.js v4 (motion) \\
Icons & Font Awesome (free) \\
Data-local & JSON/TS dataset (§10) + \texttt{localStorage} for
progress \\
Tests & Vitest + Testing Library; Playwright (E2E, offline simulate) \\
Lint/format & ESLint (flat) + Prettier (matches existing
\texttt{.pre-commit-config.yaml}) \\
Icons/app assets & inline SVG (kirik) so SW precache stays small \\
\end{longtable}

\subsubsection{12.2 Layered architecture}\label{layered-architecture}

\begin{verbatim}
┌───────────────────────────────────────────────┐
│ UI Layer        React components, routes,      │
│                 screens, story slides          │
├───────────────────────────────────────────────┤
│ Application     LevelFSM, Question session,    │
│                 progress store, reward logic   │
├───────────────────────────────────────────────┤
│ Domain          aksara.ts dataset, Stroke      │
│                 matcher interface, scoring     │
├───────────────────────────────────────────────┤
│ Infrastructure  CanvasEngine (capture+guide), │
│                 DTW matcher, Zustand persist,  │
│                 PWA/SW, audio (gamelan blips)  │
└───────────────────────────────────────────────┘
\end{verbatim}

\subsubsection{\texorpdfstring{12.3 State model (Zustand
\texttt{useProgress})}{12.3 State model (Zustand useProgress)}}\label{state-model-zustand-useprogress}

\begin{Shaded}
\begin{Highlighting}[]
\KeywordTok{interface}\NormalTok{ ProgressState \{}
\NormalTok{  completedLevels}\OperatorTok{:} \DataTypeTok{number}\NormalTok{[]}\OperatorTok{;}        \CommentTok{// [1,2,3]}
\NormalTok{  rewards}\OperatorTok{:} \DataTypeTok{string}\NormalTok{[]}\OperatorTok{;}                \CommentTok{// [\textquotesingle{}pedang\textquotesingle{},\textquotesingle{}perisai\textquotesingle{}]}
\NormalTok{  bestByQuestion}\OperatorTok{:} \BuiltInTok{Record}\OperatorTok{\textless{}}\DataTypeTok{string}\OperatorTok{,} \DataTypeTok{number}\OperatorTok{\textgreater{};} \CommentTok{// questionId {-}\textgreater{} percent score}
\NormalTok{  currentLevel}\OperatorTok{?:} \DataTypeTok{number}\OperatorTok{;}            \CommentTok{// resume point}
\NormalTok{  unfinished}\OperatorTok{:} \DataTypeTok{string}\NormalTok{[]}\OperatorTok{;}             \CommentTok{// unresolved question ids in current level}
\NormalTok{  settings}\OperatorTok{:}\NormalTok{ \{ sound}\OperatorTok{:} \DataTypeTok{boolean}\OperatorTok{;}\NormalTok{ dark}\OperatorTok{:} \DataTypeTok{boolean}\NormalTok{ \}}\OperatorTok{;}
\NormalTok{\}}
\end{Highlighting}
\end{Shaded}

\subsubsection{12.4 Offline \& SW strategy}\label{offline-sw-strategy}

\begin{itemize}
\tightlist
\item
  \textbf{Precache (build-time, hashed):} HTML shell, JS/CSS chunks,
  Javanese font (woff2), dataset, inline SVG icons --- the entire app
  works with network off.
\item
  \textbf{Runtime cache:} none required for MVP (all local); precache
  covers everything.
\item
  \textbf{Update flow:} \texttt{registerSW} + prompt-to-refresh banner
  when \texttt{updatedready}.
\item
  \textbf{No backend:} zero API calls; CSP tightened; no analytics.
\item
  Manifest: name ``Petualangan Ajisaka'', \texttt{display:\ standalone},
  \texttt{orientation:\ portrait} (writing optimised), theme color =
  \texttt{-\/-color-paper}, icon set incl.~\texttt{512px} maskable.
\end{itemize}

\begin{center}\rule{0.5\linewidth}{0.5pt}\end{center}

\subsection{13. Proposed Project
Structure}\label{proposed-project-structure}

\begin{verbatim}
javanese_learning_app/
├─ public/
│  ├─ icons/            (192, 512, maskable)
│  └─ fonts/            (NotoSansJavanese woff2)
├─ src/
│  ├─ data/             aksara.ts, levels.ts, sentences.ts
│  ├─ engine/           capture.ts, normalize.ts, dtw.ts, matcher.ts
│  ├─ state/            progress.ts
│  ├─ ui/               components/, screens/
│  ├─ hooks/            useCanvas.ts, useStrokeSession.ts
│  ├─ styles/           tokens.css, base.css
│  ├─ App.tsx
│  └─ main.tsx
├─ vite.config.ts        (+ vite-plugin-pwa)
├─ tsconfig.json
├─ index.html
└─ DESIGN_AND_DEV.md
\end{verbatim}

\begin{center}\rule{0.5\linewidth}{0.5pt}\end{center}

\subsection{14. Testing Strategy}\label{testing-strategy}

\begin{itemize}
\tightlist
\item
  \textbf{Unit:} \texttt{dtw.ts} (match/fail cases),
  \texttt{normalize.ts}, progress-store transitions, question scoring,
  level-unlock gating.
\item
  \textbf{Component:} practice screen states
  (empty/guiding/warn/error/success), canvas button labels, modal focus
  trap.
\item
  \textbf{E2E (Playwright):} flow Home→L1 completed→L2 unlocked; reload
  mid-level resumes; \textbf{offline emulation} verifies full
  playthrough with no network.
\item
  \textbf{Manual device matrix:} iPad, Android tablet (Chrome), desktop
  (portrait mobile view).
\item
  \textbf{A11y:} axe-core scan in CI.
\item
  Pre-commit (existing hooks) keep running; add
  \texttt{eslint}/\texttt{prettier}/\texttt{tsc\ -\/-noEmit} there.
\end{itemize}

\begin{center}\rule{0.5\linewidth}{0.5pt}\end{center}

\subsection{15. Performance Budget
(offline-first)}\label{performance-budget-offline-first}

\begin{longtable}[]{@{}
  >{\raggedright\arraybackslash}p{(\columnwidth - 2\tabcolsep) * \real{0.5000}}
  >{\raggedright\arraybackslash}p{(\columnwidth - 2\tabcolsep) * \real{0.5000}}@{}}
\toprule\noalign{}
\begin{minipage}[b]{\linewidth}\raggedright
Metric
\end{minipage} & \begin{minipage}[b]{\linewidth}\raggedright
Budget
\end{minipage} \\
\midrule\noalign{}
\endhead
\bottomrule\noalign{}
\endlastfoot
First load (offline cache hit) & ≤ 1.5 MB JS + assets; font ≤ 600 KB
woff2 \\
TTI on mid-range tablet & ≤ 3 s \\
SW precache total & ≤ 4--5 MB \\
Canvas frame during tracing & 60 fps \\
Recognition latency & \textless{} 50 ms (DTW on 32-pt resample) \\
\end{longtable}

\begin{center}\rule{0.5\linewidth}{0.5pt}\end{center}

\subsection{16. Accessibility (WCAG 2.1
AA)}\label{accessibility-wcag-2.1-aa}

\begin{itemize}
\tightlist
\item
  Color is never the only signal: pass/warn/error carry icons + text.
\item
  Focus-visible ring \texttt{-\/-color-focus},
  \texttt{outline-offset:\ 2px}; skip-link; landmark tags.
\item
  Keyboard operable (arrows on story slides; Enter/Space on controls);
  canvas has an accessible fallback description for each question.
\item
  \texttt{prefers-reduced-motion} disables stroke-guide animation,
  particles, slide fx.
\item
  Contrast: body text ≥ 4.5:1; UI ≥ 3:1 (verify with WebAIM checker on
  final hex).
\item
  All Indonesian copy short, plain; alt text on illustrated scenes.
\end{itemize}

\begin{center}\rule{0.5\linewidth}{0.5pt}\end{center}

\subsection{17. Security}\label{security}

\begin{itemize}
\tightlist
\item
  No user input is rendered as HTML (all content from bundled dataset).
\item
  CSP:
  \texttt{default-src\ \textquotesingle{}self\textquotesingle{};\ script-src\ \textquotesingle{}self\textquotesingle{};\ font-src\ \textquotesingle{}self\textquotesingle{};\ img-src\ \textquotesingle{}self\textquotesingle{}\ data:;}
  (no external CDNs at runtime --- everything precached).
\item
  No secrets/keys in client (none needed); no network traffic to leak.
\item
  Canvas input validated/ignored outside the write-box; multi-touch
  suppresses stray strokes.
\end{itemize}

\begin{center}\rule{0.5\linewidth}{0.5pt}\end{center}

\subsection{18. Roadmap / Milestones}\label{roadmap-milestones}

\begin{longtable}[]{@{}
  >{\raggedright\arraybackslash}p{(\columnwidth - 2\tabcolsep) * \real{0.5000}}
  >{\raggedright\arraybackslash}p{(\columnwidth - 2\tabcolsep) * \real{0.5000}}@{}}
\toprule\noalign{}
\begin{minipage}[b]{\linewidth}\raggedright
Milestone
\end{minipage} & \begin{minipage}[b]{\linewidth}\raggedright
Scope
\end{minipage} \\
\midrule\noalign{}
\endhead
\bottomrule\noalign{}
\endlastfoot
\textbf{M0 --- Seed \& shell} & Vite+TS app, tokens, routes, dashboard,
manifest/SW, empty practice screen \\
\textbf{M1 --- Stroke engine} & Canvas capture, guide animation, DTW
matcher, feedback mapping \\
\textbf{M2 --- Content} & Full dataset (Nglegena 20+ / Sandangan 8+ /
Pasangan pairs), level configs, story copy in Indonesian \\
\textbf{M3 --- Game flow} & Level FSM, rewards, Dora/villager joins,
ending, progress persistence \\
\textbf{M4 --- Polish QA} & Particles, sound, offline E2E, a11y scan,
device matrix, update banner \\
\end{longtable}

\begin{center}\rule{0.5\linewidth}{0.5pt}\end{center}

\subsection{19. Future Work (after MVP)}\label{future-work-after-mvp}

\begin{itemize}
\tightlist
\item
  Transliteration / text-to-Aksara input; pronunciation (gamelan + TTS).
\item
  Cloud sync \& teacher analytics (rename: ``Guru'' report).
\item
  Android APK via TWA; wider glyph set (Angka, Murda, Swara).
\item
  ML stroke classifier upgrade behind the matcher interface.
\end{itemize}

\begin{center}\rule{0.5\linewidth}{0.5pt}\end{center}

\subsection{20. Open Questions \&
Assumptions}\label{open-questions-assumptions}

\begin{longtable}[]{@{}
  >{\raggedright\arraybackslash}p{(\columnwidth - 4\tabcolsep) * \real{0.3333}}
  >{\raggedright\arraybackslash}p{(\columnwidth - 4\tabcolsep) * \real{0.3333}}
  >{\raggedright\arraybackslash}p{(\columnwidth - 4\tabcolsep) * \real{0.3333}}@{}}
\toprule\noalign{}
\begin{minipage}[b]{\linewidth}\raggedright
\#
\end{minipage} & \begin{minipage}[b]{\linewidth}\raggedright
Question
\end{minipage} & \begin{minipage}[b]{\linewidth}\raggedright
Proposed default / status
\end{minipage} \\
\midrule\noalign{}
\endhead
\bottomrule\noalign{}
\endlastfoot
1 & Level 1 question count & Propose \textbf{10--12} Nglegena; confirm
with curriculum. \\
2 & Exact question contents per set & Propose starter sets; finalize
with teacher input. \\
3 & Full Level-1 sinopsis paragraph & Missing from PDF extraction ---
supply copy. \\
4 & Reward/asset art source & Inline SVG illustrations; replaceable
later. \\
5 & UI framework & Recommend \textbf{React}; Svelte alternative if
bundle size is critical. \\
6 & Portrait-only PWA OK? & Yes (writing orientation); desktop shows
rotated mock. \\
\end{longtable}

\begin{center}\rule{0.5\linewidth}{0.5pt}\end{center}

\subsection{21. References}\label{references}

\begin{itemize}
\tightlist
\item
  Source narrative: \texttt{projek\ Isif.pdf} (Naskah Lengkap Aplikasi
  Petualangan Ajisaka).
\item
  Font: Noto Sans Javanese (OFL) --- self-hosted.
\item
  Standards: WCAG 2.1 AA; PWA installability criteria (Chrome/Android,
  iOS 16+).
\item
  Design system basis: playful \emph{Carnival} token set (paper/accent
  OKLCH architecture) re-hued to the keraton palette in §9.2.
\end{itemize}

\begin{center}\rule{0.5\linewidth}{0.5pt}\end{center}

\emph{End of document. Next step: confirm §20 decisions, then scaffold
M0.}\# Naskah Aplikasi Petualangan Ajisaka

\emph{Tema utama: gamifikasi keterampilan menulis Aksara Jawa.}

\begin{center}\rule{0.5\linewidth}{0.5pt}\end{center}

\subsection{Halaman Utama dan
Navigasi}\label{halaman-utama-dan-navigasi}

Aplikasi dibuka dengan halaman depan yang menampilkan judul media
pembelajaran dan satu tombol PLAY besar. Setelah tombol itu ditekan,
pemain masuk ke menu utama yang menampung lima pilihan:

\begin{enumerate}
\def\labelenumi{\arabic{enumi}.}
\tightlist
\item
  \textbf{Prolog} --- pengenalan cerita dan asal-usul.
\item
  \textbf{Level 1: Pemula} --- menulis aksara dasar.
\item
  \textbf{Level 2: Mahir} --- menulis aksara sandangan.
\item
  \textbf{Level 3: Master} --- menulis aksara pasangan dan ujian akhir.
\item
  \textbf{Kembali ke Rumah} --- kembali ke halaman depan.
\end{enumerate}

\begin{center}\rule{0.5\linewidth}{0.5pt}\end{center}

\subsection{Prolog: Mengenal Asal-Usul}\label{prolog-mengenal-asal-usul}

Pilihan pertama membawa pemain ke halaman edukasi yang membahas sejarah
terciptanya Aksara Jawa. Bagian ini menjawab dari mana aksara itu lahir
dan bagaimana ia sampai dipakai seperti sekarang, sebelum pemain mulai
menulis satu huruf pun.

\begin{center}\rule{0.5\linewidth}{0.5pt}\end{center}

\subsection{Level 1 --- Pemula: Misi di Pulau
Sanjaya}\label{level-1-pemula-misi-di-pulau-sanjaya}

Di pulau pertama, Ajisaka harus membuka segel dan mengambil pedang.
Syaratnya jelas: pemain wajib menyelesaikan tantangan praktik menulis
Aksara Jawa Dasar (Nglegena) langsung di layar digital. Proses
penulisannya dibantu penuh oleh garis panduan berbentuk anak panah, atau
\emph{stroke guide}, yang menunjukkan arah guratan tangan secara tepat
--- dari coretan pertama sampai hurufnya jadi.

Setelah tantangan menulis selesai, pemain menerima hadiah berupa
\textbf{Pedang Pusaka}. Tak lama berselang, Ajisaka bertemu seorang
pemuda lokal bernama Dora yang memutuskan menjadi pengikut setianya.
Ketika cerita selesai, muncul tombol \textbf{Next Level} untuk
melanjutkan ke menu ketiga.

\begin{center}\rule{0.5\linewidth}{0.5pt}\end{center}

\subsection{Level 2 --- Mahir: Misi di Pulau Adi
Jaya}\label{level-2-mahir-misi-di-pulau-adi-jaya}

Perjalanan berlanjut. Ajisaka dan Dora mengarungi pulau kedua, Pulau Adi
Jaya, untuk mencari pusaka berikutnya: \textbf{Perisai Sakti}. Di tengah
jalan mereka bertemu seorang warga lokal yang mengetahui tempat
penyimpanan perisai itu. Warga tersebut menantang mereka, berjanji
kembali memberi tahu lokasinya bila Ajisaka dan Dora mampu menjawab
ujian yang ia berikan.

Melihat kemahiran Ajisaka menuliskan ke-11 soal ujian itu, sang warga
kagum dan berniat ikut bertualang melawan raksasa. Ia menunjukkan letak
Perisai Sakti hingga pusaka itu berhasil diambil. Kini mereka bertiga
bersiap berlayar ke pulau utama lewat tombol \textbf{Next Level}.

\begin{center}\rule{0.5\linewidth}{0.5pt}\end{center}

\subsection{Level 3 --- Master: Pertempuran Akhir di Kerajaan
Nusantara}\label{level-3-master-pertempuran-akhir-di-kerajaan-nusantara}

Tingkat terakhir ini terbagi menjadi dua fase pertempuran menulis yang
sengit.

\subsubsection{Fase 1: Penghadangan Dua
Utusan}\label{fase-1-penghadangan-dua-utusan}

Dalam perjalanan laut menuju pusat kerajaan musuh, kapal mereka dicegat
dua utusan Raksasa Hijau yang berniat membunuh Ajisaka. Pemain wajib
menyelesaikan 20 soal Aksara Pasangan, yaitu konsonan tumpuk. Pada tiap
soal, pemain diminta menuliskan Aksara Jawa utama sekaligus bentuk
pasangannya yang menempel di bawah atau di samping. Seluruh guratan
dibimbing arah anak panah agar susunan penulisannya tidak keliru. Begitu
ke-20 soal selesai ditulis, kedua utusan itu berhasil diselesaikan.

\subsubsection{Fase 2: Ujian Penyegelan Raksasa
Hijau}\label{fase-2-ujian-penyegelan-raksasa-hijau}

Sistem langsung mengoper pemain ke menu latihan soal akhir. Ajisaka,
Dora, dan warga lokal itu berhasil menyusup ke Kerajaan Nusantara yang
sedang dikuasai Raksasa Hijau. Untuk memerdekakan kerajaan dari teror,
pemain harus menyelesaikan tiga menu latihan utama; masing-masing berisi
lima soal menulis:

\begin{enumerate}
\def\labelenumi{\arabic{enumi}.}
\tightlist
\item
  \textbf{Latihan 1} --- menulis kalimat bahasa Jawa menggunakan
  komponen Aksara Jawa Dasar.
\item
  \textbf{Latihan 2} --- menulis kalimat bahasa Jawa yang menggunakan
  aturan Aksara Sandangan.
\item
  \textbf{Latihan 3} --- menulis kalimat bahasa Jawa yang menggunakan
  aturan Aksara Pasangan.
\end{enumerate}

\subsubsection{Akhir Cerita}\label{akhir-cerita}

Setelah ketiga menu latihan penulisan, total 15 soal, diselesaikan
secara berurutan, Raksasa Hijau berhasil diselesaikan secara mutlak.
Kerajaan Nusantara kembali damai dan bebas dari penjajahan. Atas jasa
dan kepemimpinannya, pemain yang memerankan Ajisaka secara resmi
dinobatkan menjadi \textbf{Raja Kerajaan Nusantara}.\# Petualangan
Ajisaka --- Todo Pengembangan (P0--P3)

Prioritas berdasarkan \texttt{DESIGN\_AND\_DEV.md} §5 \&
\texttt{TIMELINE.md} W1--W11. \emph{Diperbarui setelah M0+M1 PoC build
(2026-08-12).}

\textbf{Skala prioritas} - \textbf{P0 --- Kritis:} tanpa ini aplikasi
tidak bisa rilis / tidak bisa dimainkan. - \textbf{P1 --- Penting:}
wajib untuk peluncuran yang baik, dapat ditunda jika terdesak. -
\textbf{P2 --- Peningkatan:} nilai tambah setelah rilis inti stabil. -
\textbf{P3 --- Masa depan:} di luar scop MVP, direncanakan kemudian.

\begin{itemize}
\tightlist
\item[$\boxtimes$]
  = selesai \& terverifikasi (build/lint/test hijau)
\item
  {[}\textasciitilde{]} = sebagian / versi PoC / butuh kurasi
\end{itemize}

\begin{center}\rule{0.5\linewidth}{0.5pt}\end{center}

\subsection{P0 --- Kritis (rintisan
rilis)}\label{p0-kritis-rintisan-rilis}

\subsubsection{Fondasi \& shell --- ✅ DONE
(M0)}\label{fondasi-shell-done-m0}

\begin{itemize}
\tightlist
\item[$\boxtimes$]
  Scaffold Vite + TypeScript strict + Tailwind v4 --- \texttt{src/}
\item[$\boxtimes$]
  Token desain keraton (OKLCH) §9.2 --- \texttt{src/src/index.css}
  (@theme)
\item[$\boxtimes$]
  Routing hash: \texttt{/}, \texttt{/menu}, \texttt{/prolog},
  \texttt{/level/1..3}, \texttt{/ending} --- App.tsx
\item[$\boxtimes$]
  Halaman Home (judul + tombol PLAY) --- FR-1.1 (Redesigned with playful
  UX)
\item[$\boxtimes$]
  Dashboard 5 menu dengan state locked/unlock --- FR-1.3, FR-1.4
  (Redesigned with Gamified LevelCards)
\item[$\boxtimes$]
  Manifest PWA + icons (192/512/maskable/apple-touch) + installable
\item[$\square$]
  Finalisasi keputusan §20 (jumlah soal L1, isi dataset, framework
  ✅React, aset ilustrasi) --- sebagian tetap terbuka
\end{itemize}

\subsubsection{Stroke engine --- PIVOT: DTW → VECTOR
GEOMETRY}\label{stroke-engine-pivot-dtw-vector-geometry}

\begin{itemize}
\tightlist
\item[$\boxtimes$]
  Hook canvas Pointer Events, HiDPI, \texttt{touch-action:\ none},
  live-ink --- useCanvasCapture
\item[$\boxtimes$]
  Normalisasi + resample (engine/geometry) --- dipakai rendering
\item[$\boxtimes$]
  Matcher \textbf{Vector Geometry (Chamfer Dist)} --- engine/raster.ts
  (Skala dinamis, cegah exploit coretan)
\item[$\boxtimes$]
  Uji matcher sintetis (square/fill/miss/offset) + 20 glyph contours ---
  vitest 8 tests
\item[$\boxtimes$]
  Umpan balik warna per guratan (green/amber/red) --- §11.4
\item[$\boxtimes$]
  Kontrol Clear (``Bersihkan'') --- FR-3.6
\item[$\boxtimes$]
  Kalibrasi toleransi coverage (COVERAGE\_PASS/WARN) dgn uji tablet
  nyata --- sudah dinamis per luas/tebal glyph
\end{itemize}

\subsubsection{Konten}\label{konten}

\begin{itemize}
\tightlist
\item[$\boxtimes$]
  Dataset Nglegena \textbf{20 contour} (dari Noto Sans Javanese,
  Y-corrected, 80-pt) --- src/data/nglegena\_contours.json
\item[$\boxtimes$]
  Dataset Sandangan 8+ core + kombinasi --- §10.1 (Label \& Dotted
  Circle diperbaiki)
\item[$\boxtimes$]
  Dataset Pasangan (20 pasang) --- §10.1 (Berhasil diekstrak ulang penuh
  dengan loops via FontTools Python script)
\item[$\boxtimes$]
  Level config L1--L3 bukti: starter ha/na/ca/ra/ka; soal 11/20/3×5
  belum terisi penuh --- §10.2 (Dataset penuh digunakan untuk L1-L3)
\item[$\boxtimes$]
  Copy cerita Indonesia penuh (dari \texttt{NASKAH\_AJISAKA.md}) ke tiap
  level --- Selesai via i18n layer
\item[$\boxtimes$]
  Precache font Javanese (woff2, 89KB, self-host) --- FR-6.2
\item[$\square$]
  \textbf{Stroke data kurasi internal (guide pen-stroke ori)} --- kanvas
  guide saat ini solid-fill contour; untuk guratan internal/order
  diajarkan butuh data kurator (§20)
\end{itemize}

\subsubsection{Game flow \& state}\label{game-flow-state}

\begin{itemize}
\tightlist
\item[$\boxtimes$]
  Level FSM path: intro screen → practice → done page
\item[$\boxtimes$]
  Reward Pedang Pusaka (L1) / Perisai Sakti (L2) / Raja (L3) ---
  LevelDone
\item[$\boxtimes$]
  Event ``Dora bergabung'' \& ``warga bergabung'' (copy) --- LevelDone
\item[$\boxtimes$]
  Unlock chain + tombol Next Level --- FR-4.3
\item[$\boxtimes$]
  Interstitial Fase 1 / Fase 2 pilihan menu non-linear di Level 3 ---
  FR-4.4 (Selesai, tracking granular per phase)
\item[$\boxtimes$]
  Ending dinobatkan Raja terhubung setelah L3 Fase 1 \& 2 selesai ---
  Selesai
\item[$\boxtimes$]
  Persist progress (Zustand persist): completedLevels, completedPhases,
  rewards --- FR-5.1
\item[$\square$]
  Resume level saat reload --- FR-5.2 (progress tersimpan, tapi resume
  per-pertanyaan di tengah jalan belum)
\item[$\square$]
  Latihan soal L1: mulai dari soal 1 → reward→ level 2 dsb (starter set
  dipakai ulang utk L2/L3 sementara)
\end{itemize}

\subsubsection{Fitur Tambahan \& UX Polish --- ✅
DONE}\label{fitur-tambahan-ux-polish-done}

\begin{itemize}
\tightlist
\item[$\boxtimes$]
  Localization penuh: Indonesia, English, Basa Jawa (Krama Inggil) via
  i18next + ekstraksi data layer
\item[$\boxtimes$]
  Audio engine Gamelan Web Audio API (Ombak, saron, gong ageng, keprak,
  detuned partials)
\item[$\boxtimes$]
  Hide global scrollbars untuk PWA fullscreen feel
\item[$\boxtimes$]
  Mengunci Orientasi Layar (Portrait PWA Lock)
\end{itemize}

\subsubsection{PWA \& offline}\label{pwa-offline}

\begin{itemize}
\tightlist
\item[$\boxtimes$]
  Service worker (Workbox) precache penuh (\textasciitilde500 KB, 14
  entries) --- FR-6.2
\item[$\boxtimes$]
  Update banner prompt-refresh --- FR-6.3 (registerSW onNeedRefresh,
  diperbarui ke autoUpdate)
\item[$\boxtimes$]
  Uji playthrough offline penuh di perangkat (IOS icons +
  user-scalable=no)
\end{itemize}

\subsubsection{QA rilis}\label{qa-rilis}

\begin{itemize}
\tightlist
\item[$\boxtimes$]
  Unit test matcher/raster + geometry + dtw (8 tests) --- vitest
\item[$\square$]
  Coverage ≥90\% utk engine (target paket tambahan)
\item[$\square$]
  E2E Playwright: Home→L1 selesai→L2 unlock; reload mid-level
\item[$\square$]
  Matrix perangkat (iPad, Android tablet, desktop) --- §14
\end{itemize}

\begin{center}\rule{0.5\linewidth}{0.5pt}\end{center}

\subsection{P1 --- Penting (peluncuran
berkualitas)}\label{p1-penting-peluncuran-berkualitas}

\begin{itemize}
\tightlist
\item[$\boxtimes$]
  Partikel sukses per soal \& reward (≤60 particle) --- §9.5 (Selesai
  via Anime.js confetti)
\item[$\boxtimes$]
  Suara gamelan ringan (Web Audio) (engine audio.ts selesai) --- §12.2
\item[$\boxtimes$]
  Transisi antarmuka \& animasi UI (menggunakan animejs dgn
  bounce/stagger)
\item[$\boxtimes$]
  Transisi slide sinopsis (Prolog.tsx) --- FR-2.2
\item[$\square$]
  Modal congrats/fokus trap --- §9.7
\item[$\boxtimes$]
  Audit a11y axe (P0 kritikal/serius = 0, kontras AA) --- §16
\item[$\boxtimes$]
  Skip-link + landmark announcement + \texttt{:focus-visible} penuh ---
  §16
\item[$\boxtimes$]
  Suport keyboard (story slides, Enter/Space) --- §16 (Selesai untuk
  Prolog, Practice, Phase2, LevelDone)
\item[$\square$]
  Alternatif aksesibel untuk canvas (deskripsi soal lengkap) --- §16
\item[$\square$]
  Optimasi bundle ≤1,5 MB + font ≤600 KB (saat ini ±500 KB total
  precache) --- §15
\item[$\square$]
  Pola 60fps saat menulis dgn guide menyala (rAF + dpr scaling)--- §15
\item
  {[}\textasciitilde{]} Kanvas toggle panduan (Tampilkan/Sembunyikan
  Contoh) --- selesai, verifikasi UX
\item[$\square$]
  E2E offline emulation penuh (Playwright)
\end{itemize}

\begin{center}\rule{0.5\linewidth}{0.5pt}\end{center}

\subsection{P2 --- Peningkatan (setelah rilis inti
stabil)}\label{p2-peningkatan-setelah-rilis-inti-stabil}

\begin{itemize}
\tightlist
\item[$\square$]
  TTS pengucapan fonem + pelafalan nama aksara (lokal, gabung dgn suara)
\item[$\boxtimes$]
  Transliterasi teks → Aksara Jawa (masukan bebas) --- Selesai via
  ``Free Type''
\item[$\square$]
  Set glyph lanjutan: Angka, Murda, Swara
\item[$\square$]
  Peningkatan model matcher → classifier ML (ONNX/TFLite \textless500
  KB) --- §11.3
\item[$\square$]
  Statistik per latihan tambahan (skor, waktu, retry)
\item[$\square$]
  Kuis pengetahuan cepat (bukan hanya menulis) sbg penguat
\item[$\square$]
  Mode bantuan ekstra anak kecil (petunjuk lebih besar, warna kontras
  tinggi)
\item[$\square$]
  Pengaturan lanjutan (volume, mode kontras tinggi, ukuran teks)
\end{itemize}

\begin{center}\rule{0.5\linewidth}{0.5pt}\end{center}

\subsection{P3 --- Masa depan (di luar
MVP)}\label{p3-masa-depan-di-luar-mvp}

\begin{itemize}
\tightlist
\item[$\boxtimes$]
  Dockerization untuk hosting production yang mandiri
\item[$\square$]
  Cloud sync \& laporan guru (``Guru Report'') --- §19
\item[$\square$]
  Leaderboard lokal/kelas + akun opsional --- G5, §19
\item[$\square$]
  APK Android via TWA; wrapper iPad --- §19
\item[$\boxtimes$]
  Multi-bahasa (Jawa--Indonesia--English) --- Selesai (i18n ditambahkan
  ke MVP)
\item[$\square$]
  Editor stroke konten untuk guru (impor data aksara baru)
\item[$\square$]
  Ekspor sertifikat penyelesaian (printer-friendly)
\end{itemize}

\begin{center}\rule{0.5\linewidth}{0.5pt}\end{center}

\subsection{Notulensi Pengunci
Keputusan}\label{notulensi-pengunci-keputusan}

\begin{longtable}[]{@{}
  >{\raggedright\arraybackslash}p{(\columnwidth - 6\tabcolsep) * \real{0.2500}}
  >{\raggedright\arraybackslash}p{(\columnwidth - 6\tabcolsep) * \real{0.2500}}
  >{\raggedright\arraybackslash}p{(\columnwidth - 6\tabcolsep) * \real{0.2500}}
  >{\raggedright\arraybackslash}p{(\columnwidth - 6\tabcolsep) * \real{0.2500}}@{}}
\toprule\noalign{}
\begin{minipage}[b]{\linewidth}\raggedright
\#
\end{minipage} & \begin{minipage}[b]{\linewidth}\raggedright
Keputusan
\end{minipage} & \begin{minipage}[b]{\linewidth}\raggedright
Status
\end{minipage} & \begin{minipage}[b]{\linewidth}\raggedright
Deadlines
\end{minipage} \\
\midrule\noalign{}
\endhead
\bottomrule\noalign{}
\endlastfoot
1 & Jumlah soal Level 1 (starter: 5 glyph ha/na/ca/ra/ka) & ⏳ terbuka &
sebelum isi dataset penuh \\
2 & Isi dataset (glyph mana saja per level) & ⏳ terbuka & sebelum
dataset lengkap \\
3 & Sinopsis L1 (copy lengkap dari PDF hilang) & ⏳ terbuka & sebelum
copy cerita penuh \\
4 & Framework React vs Svelte & ✅ \textbf{React} (diputuskan) & --- \\
5 & Aset ilustrasi (SVG inline) & ⏳ terbuka & akhir minggu 2 \\
6 & Model evaluasi menulis & ✅ \textbf{Raster coverage} (DTW outline
diganti) & --- \\
7 & Sumber data pen-stroke kurasi (guide internal) & ⏳ terbuka & P1 \\
\end{longtable}

\begin{center}\rule{0.5\linewidth}{0.5pt}\end{center}

\emph{Perbarui notulensi § berdasarkan keputusan TIMELINE.md; kotak cek
dicoret saat selesai.}

\end{document}
